\chapter*{前言}
\addcontentsline{toc}{chapter}{前言}

操作系统不是一组系统调用名，也不是“进程、内存、文件、网络”这些名词的并列清单。它是一套把硬件资源变成可控抽象的工程协议：CPU 时间被包装成可调度实体，物理内存被包装成带权限的地址空间，设备被包装成文件和事件，崩溃风险被包装成日志、提交点和恢复规则。

本册是一部操作系统原理教材。它从一台没有操作系统的机器开始，但不重复数字电路、处理器数据通路和总线结构；这些内容由独立硬件卷完整讲解。本册只取已经成立的硬件事实，逐步追问：一项程序怎样装入，协调代码怎样常驻，处理器怎样从不合作的程序手中重新取得控制，多项执行怎样拥有可保存的身份，以及受限程序怎样请求内核服务。

正文以一套从复位向量开始的 x86-64 教学系统为主体架构。项目当前已经完成自研
ROM、Stage 1、长模式、ELF64 内核交接、异常、物理页、四级页表、早期堆以及
传统中断与设备闭环；用户边界、进程调度、同步、文件系统和用户环境是后续
里程碑。每章先沿项目解决一项具体失败，再把该最小实现扩展到多任务、多核、
并发和故障条件，形成任务、系统调用、地址空间、页缓存、设备请求、VFS、
文件系统、网络、安全、恢复与观测等完整机制。

当这些对象都已经获得明确定义，书末再沿一次文件保存请求完成全书综合：执行流如何被调度，用户请求如何进入内核，地址如何验证，文件对象如何引用，设备请求如何异步完成，数据又如何跨越持久化边界。此时全景不是术语清单，而是读者能够逐段核对的因果链。

本册通常先给出具体问题或反例，再分析直接方案的局限，继而建立相应机制、不变量与观察证据。读者不应只记住“进程是资源分配单位，线程是调度单位”这样的概括，而应能够回答：旧机器为什么失败，新增对象保存哪些状态，谁有权改变状态，错误路径如何退出，哪些证据能够支持结论。

本册仍以原理分析为主，但项目不再只是章末例子。复位地址、磁盘描述符、处理器
模式、BootInfo、异常帧、页帧状态和 IRQ 确认等真实对象构成章节之间的连续机器
状态；现有项目尚未达到的高级机制会明确标成设计里程碑，而不会冒充已完成实现。

\begin{keyidea}
操作系统的每个特性都应该能回答三个问题：它依赖哪个硬件事实，它修复了前一阶段的哪个失败，又引入了哪些新的状态、成本和测试责任。
\end{keyidea}
